\begin{textsample}{POS Dim 3 – ap – Score 15.00 – ap/ap\_inf\_12.txt}  \label{ex:f3_pos_141}
IV – Campo de Experiência : \textbf{Escuta}, Fala, Pensamento e Imaginação : Desde o nascimento, as \textbf{crianças} participam de situações comunicativas cotidianas com as pessoas com as quais interagem.
As primeiras formas de interação do \textbf{bebê} são os movimentos do seu corpo, o olhar, a postura corporal, o sorriso, o choro e outros recursos vocais, que ganham sentido com a interpretação do outro. \textbf{Progressivamente}, as \textbf{crianças} vão ampliando e enriquecendo seu vocabulário e demais recursos de expressão e de compreensão, apropriando -se da língua materna – que se torna, pouco a pouco, seu veículo privilegiado de interação.
Na Educação \textbf{Infantil}, é importante promover experiências nas quais as \textbf{crianças} possam falar e ouvir, potencializando sua participação na cultura oral, pois é na \textbf{escuta} de histórias, na participação em conversas, nas descrições, nas narrativas elaboradas individualmente ou em grupo e nas implicações com as múltiplas linguagens que a \textbf{criança} se constitui ativamente como sujeito singular e pertencente a um grupo social.
Desde cedo, a \textbf{criança} manifesta \textbf{curiosidade} com relação à cultura escrita : ao ouvir e acompanhar a leitura de textos, ao observar os muitos textos que circulam no contexto familiar, comunitário e escolar, ela vai construindo sua concepção de língua escrita, reconhecendo diferentes usos sociais da escrita, dos gêneros, suportes e portadores.
Na Educação \textbf{Infantil}, a imersão na cultura escrita deve partir do que as \textbf{crianças} conhecem e das \textbf{curiosidades} que deixam transparecer.
As experiências com a literatura \textbf{infantil}, propostas pelo educador, mediador entre os textos e as \textbf{crianças}, contribuem para o desenvolvimento do gosto pela leitura, do estímulo à imaginação e da ampliação do conhecimento de mundo.
Além disso, o contato com histórias, contos, fábulas, poemas, cordéis etc. propicia a familiaridade com livros, com diferentes gêneros literários, a diferenciação entre ilustrações e escrita, a aprendizagem da direção da escrita e as formas corretas de manipulação de livros.
Nesse convívio com textos escritos, as \textbf{crianças} vão construindo hipóteses sobre a escrita que se revelam, inicialmente, em rabiscos e garatujas e, à medida que vão conhecendo letras, em escritas espontâneas, não convencionais, mas já indicativas da compreensão da escrita como sistema de representação da língua.
Direitos de aprendizagem com foco no campo de experiência \textbf{Escuta}, fala, pensamento e imaginação : - CONVIVER com \textbf{crianças} e \textbf{adultos}, compartilhando situações comunicativas cotidianas, constituindo modos de pensar, imaginar, sentir, narrar, dialogar e conhecer. - \textbf{BRINCAR} com parlendas, trava-línguas, adivinhas, textos de memória, ( \textbf{brincadeiras} de ) rodas, \textbf{brincadeiras} cantadas ( \textbf{brinquedos} cantados ) e jogos, ampliando o repertório das manifestações culturais da tradição local e de outras culturas, enriquecendo a linguagem oral, corporal, musical, \textbf{dramática}, escrita, dentre outras. - PARTICIPAR de rodas de conversa, de relatos de experiências, de contação e leitura de histórias e poesias, de construção de narrativas, da elaboração e descrição de papéis no faz de conta, da exploração de materiais impressos, analisando as estratégias comunicativas, as variedades linguísticas e descobrindo as diversas formas de organizar o pensamento. - EXPLORAR \textbf{gestos}, expressões, \textbf{sons} da língua, rimas, imagens, textos escritos, além dos sentidos das falas cotidianas, das palavras nas poesias, parlendas, canções e nos enredos de histórias, apropriando -se desses elementos para criar novas falas, enredos, histórias e escritas, convencionais ou não. - EXPRESSAR sentimentos, ideias, percepções, \textbf{desejos}, necessidades, pontos de vista, informações, dúvidas e descobertas, utilizando múltiplas linguagens, entendendo e considerando o que é comunicado pelos \textbf{colegas} e \textbf{adultos}. - CONHECER -SE, a partir de uma apropriação autoral da(s) linguagens (linguagem(ens)), interagindo com os outros, reconhecendo suas preferências por pessoas, \textbf{brincadeiras}, lugares, ( e ) histórias.

% matched lemmas: adulto, bebê, brincadeira, brincar, brinquedo, colega, criança, curiosidade, desejo, dramático, escuta, gesto, infantil, progressivamente, som
\end{textsample}
